%==============================================================================
%  mammoviz-tda — Acceso a datasets longitudinales de mamografia
%  Documento independiente de la memoria del TFG (tfg/main.tex)
%==============================================================================
\documentclass[11pt,a4paper]{article}

\usepackage[utf8]{inputenc}
\usepackage[T1]{fontenc}
\usepackage[spanish,es-tabla]{babel}
\usepackage{lmodern}

\usepackage[a4paper,top=2.5cm,bottom=2.5cm,left=2.8cm,right=2.8cm]{geometry}
\usepackage{setspace}
\onehalfspacing
\usepackage{emptypage}

\usepackage{graphicx}
\usepackage{float}
\usepackage{booktabs}
\usepackage{array}
\usepackage{caption}
\usepackage{xcolor}
\definecolor{bg}{rgb}{0.96,0.96,0.96}
\definecolor{accent}{rgb}{0.11,0.35,0.60}

\usepackage{listings}
\lstset{
  basicstyle=\ttfamily\small,
  backgroundcolor=\color{bg},
  keywordstyle=\color{blue!70!black}\bfseries,
  commentstyle=\color{green!40!black},
  stringstyle=\color{red!60!black},
  frame=single, breaklines=true, showstringspaces=false,
  captionpos=b, language=bash, columns=flexible,
}

\usepackage[hidelinks,colorlinks=true,linkcolor=accent,urlcolor=accent]{hyperref}
\usepackage{enumitem}
\usepackage{fancyhdr}
\pagestyle{fancy}
\fancyhf{}
\fancyhead[L]{mammoviz-tda — acceso a datasets longitudinales}
\fancyhead[R]{\thepage}
\renewcommand{\headrulewidth}{0.4pt}

\title{\textbf{mammoviz-tda}\\[4pt]\large Acceso a datasets longitudinales de mamografía\\[2pt]\normalsize (CSAW-CC, EMBED y alternativas investigadas)}
\author{}
\date{Noviembre 2026}

\begin{document}
\maketitle
\thispagestyle{empty}

\begin{abstract}
\noindent Este documento recoge, para no tener que volver a investigarlo,
los pasos concretos para solicitar acceso a los dos datasets longitudinales
de mamografía identificados como candidatos para evaluar la hipótesis
principal del TFG (issue \#92 del repositorio), y las demás alternativas
que se han investigado y por qué no resuelven el problema mejor que estas
dos. Verificado el 9 de noviembre de 2026 contra las páginas oficiales de
cada organización (no son suposiciones ni resúmenes de terceros sin
contrastar).
\end{abstract}

\section{Resumen rápido}
\begin{itemize}
  \item No existe ningún dataset longitudinal de mamografía de descarga
        directa sin solicitud. Los dos mejores candidatos siguen siendo
        \textbf{CSAW-CC} y \textbf{EMBED}; ninguna alternativa investigada
        es más fácil o más rápida.
  \item Recomendación: \textbf{solicitar acceso a ambos en paralelo} desde
        ahora, porque ninguno de los dos portales indica un plazo de
        respuesta, y no hay coste por solicitar a los dos.
  \item El dataset de Kaggle/RSNA (2023) es de descarga libre e inmediata,
        pero \textbf{no es longitudinal} (una sola exploración por caso):
        solo sirve como dataset adicional para la clasificación aislada
        (O3--O6 del TFG), no para la hipótesis principal (O7--O8).
\end{itemize}

\section{CSAW-CC (Karolinska Institutet)}
\label{sec:csawcc}
Gestionado a través del sistema \textbf{DORIS} de SND (Swedish National
Data Service), no por la propia Karolinska.

\begin{enumerate}[leftmargin=1.2cm]
  \item Ficha del dataset:
        \url{https://researchdata.se/catalogue/dataset/2021-204-1/1}
        (DOI \texttt{10.5878/45vm-t798}). La licencia aparece como
        \texttt{NON\_PUBLIC}/\texttt{restrictedAccess}: \emph{``available
        upon request due to the size of the material''}.
  \item Formulario de solicitud (no requiere login):
        \url{https://snd.se/en/research-data-support/data-access-requests}
  \item Datos a rellenar: tus datos de contacto y una descripción del uso
        que le vas a dar. \textbf{Importante siendo estudiante}: el
        formulario también pide el nombre y email de tu tutor/director de
        TFG, a quien SND notifica para verificar la solicitud.
  \item SND reenvía la solicitud a la organización responsable
        (Karolinska Institutet), que decide si te da acceso. No se indica
        plazo de respuesta en ninguna página oficial.
  \item Contactos directos si el formulario falla o no hay respuesta:
        \begin{itemize}
          \item Fredrik Strand (creador del dataset, Karolinska Institutet)
                --- \texttt{fredrik.strand@ki.se}
          \item Kevin Smith (KTH / SciLifeLab) --- \texttt{ksmith@kth.se}
        \end{itemize}
\end{enumerate}

\subsection*{Nota sobre el CSAW original (no CC)}
Existe un dataset "hermano" más grande, el CSAW original (499\,807 mujeres,
2008--2015), del mismo grupo. Una página de un vendedor de datos de
terceros afirma que un 20\,\% de CSAW está disponible en Kaggle sin
solicitud, \textbf{pero esta afirmación no se ha podido confirmar contra
ninguna fuente primaria} (ni la ficha de SND, ni Karolinska, ni una página
oficial de Kaggle). No usar esta cifra sin verificarla de nuevo antes de
dar por hecho que hay un atajo.

\section{EMBED (Emory University)}
\label{sec:embed}
Gestionado por el propio equipo de Emory (HITI Lab) a través de un portal
dedicado, \texttt{data.hitilab.com}.

\begin{enumerate}[leftmargin=1.2cm]
  \item Portal de solicitud: \url{https://data.hitilab.com/} $\to$ botón
        ``Request access'' (\texttt{/sign-up}).
  \item Proceso (confirmado en la propia web):
        \begin{enumerate}
          \item Enviar tu \textbf{propósito de investigación} y tu
                \textbf{afiliación institucional} para revisión.
          \item Si se aprueba, recibes credenciales de AWS \textbf{temporales},
                con alcance limitado a ese dataset concreto.
          \item Descargas directamente del bucket S3 privado con esas
                credenciales.
        \end{enumerate}
  \item Acceso limitado a investigadores académicos; usa tu correo
        institucional, no uno personal, al solicitar.
  \item Licencia / acuerdo de uso de datos (revisar antes de solicitar):
        \url{https://github.com/Emory-HITI/EMBED_Open_Data/blob/main/EMBED_license.md}
  \item Documentación general del dataset:
        \url{https://github.com/Emory-HITI/EMBED_Open_Data}
  \item Contacto directo de respaldo: Hari Trivedi ---
        \texttt{hari.trivedi@emory.edu}
\end{enumerate}

\subsection*{Aclaración importante sobre el ``20\,\% público''}
El bucket de AWS Open Data de EMBED (\texttt{s3://embed-dataset-open/})
está marcado en el registro oficial de AWS como \textbf{``Controlled
Access''}, no de acceso público. El ``20\,\%'' que aparece en varias
páginas (incluida la propia ficha de AWS) se refiere al \textbf{tamaño
relativo} de ese subconjunto frente al dataset completo (3.4M imágenes),
\textbf{no} a que sea descargable sin pedir acceso. Se comprobó
empíricamente (petición HTTP directa y \texttt{boto3} sin firmar) que el
listado anónimo del bucket devuelve \texttt{AccessDenied}. Hay que pasar
por el mismo portal \texttt{data.hitilab.com} tanto para el subconjunto
``20\,\%'' como para el dataset completo.

\section{Alternativas investigadas (y por qué no resuelven el problema)}
\label{sec:alternativas}

\begin{description}[leftmargin=0cm,style=nextline]
  \item[RSNA Screening Mammography Breast Cancer Detection (Kaggle, 2023)]
        Descarga \textbf{libre e inmediata}, sin solicitud ---
        \url{https://www.rsna.org/education/ai-resources-and-training/ai-image-challenge/screening-mammography-breast-cancer-detection-ai-challenge}.
        Pero son \textbf{exploraciones puntuales} (screening de un solo
        momento por caso, con solo 2.1\,\% de casos positivos); no incluye
        pares de exploraciones sucesivas de la misma paciente con cambio de
        diagnóstico. Útil como dataset adicional para la clasificación
        aislada (base del TFG, O3--O6), no para O7--O8.

  \item[OPTIMAM / OMI-DB (Cancer Research UK + Royal Surrey NHS)]
        Es genuinamente longitudinal (cribados seriados durante 10 años,
        172\,282 mujeres a mayo de 2020) y por tanto el candidato más
        parecido en tamaño/calidad a CSAW-CC/EMBED. Pero el acceso es
        \textbf{más burocrático, no menos}: se gestiona como oportunidad de
        licencia comercial a través de \textbf{Cancer Research Horizons}
        (\url{https://www.cancerresearchhorizons.com/licensing-opportunities/optimam-mammography-image-database-omi-db}),
        con aprobación de un comité de acceso a datos y un acuerdo de
        compartición de datos (\textit{data sharing agreement}) por
        proyecto. No parece razonable para el plazo de un TFG salvo que el
        tutor tenga ya una relación institucional con Cancer Research UK.

  \item[NYU Breast Cancer Screening Dataset (NYU Langone)]
        \textbf{Prácticamente cerrado}: es una base de datos privada de NYU
        Langone que, según sus propios autores, \emph{no se puede hacer
        pública}. Solo ofrecen evaluar modelos de otros grupos sobre su
        partición de test, previa solicitud a
        \texttt{k.j.geras@nyu.edu} (respuesta en una semana, según un
        paper reciente que usa este dato, no verificado en una página
        oficial). No es una vía de acceso a las imágenes en bruto. Como
        alternativa parcial, el modelo ya entrenado por este mismo grupo
        (Wu et al. 2019) es público en
        \url{https://github.com/nyukat/breast_cancer_classifier}: no sirve
        para evaluar la hipótesis topológica, pero podría ser un baseline
        de comparación de \textit{deep learning} sin necesitar sus datos.

  \item[Otros datasets de mamografía de descarga libre (INbreast, VinDr-Mammo,
        CMMD, MIAS)]
        Todos son de \textbf{una sola exploración por caso}; ninguno tiene
        estructura longitudinal (exploraciones repetidas de la misma
        paciente en distintas fechas). No son sustitutos para O7--O8.
\end{description}

\section{Conclusión}
Tras esta revisión, \textbf{CSAW-CC y EMBED siguen siendo las dos mejores
(y únicas razonables) opciones} para evaluar la hipótesis de detección
anticipada sobre datos reales. No hay ningún atajo de descarga directa que
se haya podido confirmar. La recomendación práctica es solicitar acceso a
ambos ya (Secciones~\ref{sec:csawcc} y~\ref{sec:embed}), en paralelo, dado
que:
\begin{itemize}
  \item Ninguno de los dos indica un plazo de respuesta, así que cuanto
        antes se solicite, antes se puede empezar O8 en cuanto llegue la
        aprobación.
  \item No hay ningún coste ni exclusividad en solicitar a los dos a la
        vez; el que responda primero (o dé mejores condiciones) se usa
        para O8, y el otro queda como respaldo o para una comparación
        cruzada si ambos llegaran a concederse.
  \item Mientras se espera, el trabajo puede seguir avanzando con la
        validación sintética ya hecha (\texttt{tfg/capitulos/05-metodologia.tex},
        sección ``Metodología longitudinal'') y con la ampliación del
        piloto real de CBIS-DDSM para la base de clasificación (issue
        \#82), que no dependen de este trámite.
\end{itemize}

\end{document}
